\chapter*{Resume}
\addcontentsline{toc}{chapter}{Resume}

Les organisations exposees a Internet accumulent des milliers de vulnerabilites
detectees par des scanners tels que Qualys, sans toujours disposer d'un moyen fiable
de les hierarchiser selon le risque reel qu'elles font courir au systeme d'information.
Une severite technique brute (CVSS ou QDS) ne suffit pas : un meme defaut n'a pas le
meme impact selon qu'il touche un serveur de test isole ou une base de donnees de
production exposee sur Internet. Ce projet, realise dans le cadre d'un stage de fin
d'annee au sein de \organisme, propose une plateforme qui
croise automatiquement les exports Qualys (vulnerabilites) et CMDB (criticite metier
des actifs) pour calculer un score de risque contextualise par actif, l'\emph{Asset
Risk Score} (ARS), et pour en deriver une gouvernance operationnelle (classes de
risque, delais de remediation SLA, mesure de la dette operationnelle via le MTTR
pondere).

Le travail effectue couvre l'ensemble de la chaine applicative : l'ingestion
asynchrone et par lots (Spring Batch) des fichiers XML Qualys et CMDB avec
dedoublonnage par empreinte cryptographique, la modelisation mathematique et
l'implementation du moteur de calcul de risque, la construction d'un tableau de bord
temps reel (Server-Sent Events, Chart.js) restituant la posture de securite du parc
sous forme de graphiques (repartition par classe, diagramme de Sankey, historique,
conformite SLA), ainsi qu'un assistant conversationnel fonde sur une architecture RAG
(Retrieval-Augmented Generation) permettant d'interroger ces donnees en langage
naturel -- ce dernier volet restant, au moment de la redaction, au stade du
developpement et de la conception (prototype fonctionnel non encore valide pour un
environnement de recette ou de production). Ce rapport detaille le contexte, la conception, les fondements
mathematiques -- demontres formellement -- du modele de risque, la realisation
technique et les perspectives d'evolution de la plateforme.

\vspace{0.5cm}
\noindent\textbf{Mots-cles :} gestion des vulnerabilites, risque cyber, CMDB, Qualys,
Asset Risk Score, MTTR, Spring Batch, tableau de bord temps reel, RAG, Spring Boot.

\bigskip
\bigskip

\chapter*{Abstract}
\addcontentsline{toc}{chapter}{Abstract}

Organisations exposed to the Internet accumulate thousands of vulnerabilities
detected by scanners such as Qualys, without always having a reliable way to
prioritise them according to the actual risk they pose to the information system. A
raw technical severity score (CVSS or QDS) is not sufficient: the same flaw does not
carry the same impact depending on whether it affects an isolated test server or an
Internet-facing production database. This project, carried out as part of a final
year internship at \organisme, introduces a platform that
automatically cross-references Qualys exports (vulnerabilities) with CMDB exports
(business criticality of assets) to compute a contextualised risk score per asset,
the \emph{Asset Risk Score} (ARS), and to derive an operational governance layer from
it (risk classes, SLA remediation deadlines, operational debt measurement through a
weighted MTTR).

The work covers the whole application chain: asynchronous, batch-oriented ingestion
(Spring Batch) of Qualys and CMDB XML files with deduplication by cryptographic hash,
the mathematical modelling and implementation of the risk-scoring engine, a
real-time dashboard (Server-Sent Events, Chart.js) rendering the fleet's security
posture through charts (risk-class distribution, Sankey diagram, historical trend,
SLA compliance), and a conversational assistant built on a Retrieval-Augmented
Generation (RAG) architecture allowing this data to be queried in natural language --
this last component still being, at the time of writing, at the development and
design stage (a functional prototype not yet validated for a staging or production
environment). This report details the context, the design, the mathematically proven foundations
of the risk model, the technical implementation and the future perspectives of the
platform.

\vspace{0.5cm}
\noindent\textbf{Keywords:} vulnerability management, cyber risk, CMDB, Qualys, Asset
Risk Score, MTTR, Spring Batch, real-time dashboard, RAG, Spring Boot.

\clearpage
